\setcounter{chaptercntr}{3}

\sectionbreak {
\gostTitleFont
\redline
\thechaptercntr \space
РЕАЛИЗАЦИЯ ОСНОВНЫХ МОДУЛЕЙ СИСТЕМЫ
}

\titlespace

\subsection*{
  \gostTitleFont
  \redline
  \thechaptercntr .\thesubchaptercntr \spc
  Реализация сервиса управления заказами
} \addtocounter{subchaptercntr}{1}

\subtitlespace

{\gostFont

  \par \redline Сервис поездок является центральным компонентом системы, отвечающим за весь жизненный цикл заказа - от момента создания до выставления оценки. Именно в нём сосредоточена основная бизнес-логика: создание поездки, расчёт стоимости, назначение водителя и перевод заказа между состояниями. Сервис реализован на языке Java с использованием Spring Boot 3.4 и Spring Data JPA. Данные хранятся в отдельной БД ride\_db, для управления схемой которой применяется Liquibase [3,6].

  \par \redline Проект сервиса организован по стандартной слоистой структуре. Слой контроллеров принимает HTTP-запросы и возвращает ответы в формате JSON. Слой сервисов содержит бизнес-логику - именно в нём реализованы методы создания поездки, назначения водителя и перевода между статусами. Слой доступа к данным построен на Spring Data JPA и включает репозитории для сущностей Ride, Address и KafkaMessage. Для преобразования между сущностями и объектами передачи данных используются мапперы на базе MapStruct. Такая структура позволяет чётко разделить ответственность между компонентами и упрощает тестирование [3,4].

  \par \redline В основе сервиса лежат две основные сущности JPA: Ride и Address. Ride описывает поездку и связан с двумя записями Address через внешние ключи departure\_address и destination\_address. Address хранит геоинформацию и используется повторно: если один и тот же адрес указывается в разных заказах, новая запись в таблице address не создаётся, а используется существующая. В таблице 4.1 приведены основные поля сущности Ride.

  \topTablespace
  {\begin{Center}
      \par Таблица \thechaptercntr .\thetablecntr \spc {--} Основные поля сущности Ride

      \begin{tabular}{|c|c|c|}
        \hline
        Поле & Тип & Описание \\ \hline
        ride\_id & Long & Первичный ключ, автогенерация \\ \hline
        departure\_address\_id & Long (FK) & Ссылка на адрес отправления \\ \hline
        destination\_address\_id & Long (FK) & Ссылка на адрес назначения \\ \hline
        distance & Double & Дистанция поездки в метрах \\ \hline
        driver\_id & Long & Идентификатор водителя (логический FK) \\ \hline
        passenger\_id & Long & Идентификатор пассажира (логический FK) \\ \hline
        ride\_status & Integer & Статус поездки, значение от 0 до 11 \\ \hline
        ride\_price & BigDecimal & Стоимость в условных единицах \\ \hline
        pickup\_time & LocalDateTime & Время подачи автомобиля \\ \hline
        complete\_time & LocalDateTime & Время завершения поездки \\ \hline
        order\_create\_time & LocalDateTime & Время создания заказа \\ \hline
    \end{tabular} \end{Center}}
  \botTablespace \addtocounter{tablecntr}{1}

  \par \redline Статус поездки хранится в БД в виде целого числа, а в коде преобразуется в перечисление RideStatus из двенадцати значений. Начальное состояние - CREATED. Затем заказ последовательно проходит цепочку WAITING\_FOR\_DRIVER, DRIVER\_ASSIGNED, ACCEPTED, ON\_THE\_WAY\_TO\_PASSENGER, PASSENGER\_PICKED\_UP, ON\_THE\_WAY, COMPLETED, PAYMENT\_PENDING, PAID и RATE. Из любого состояния до COMPLETED предусмотрен переход в CANCELLED. Преобразование между числом и перечислением выполняет конвертер RideStatusConverter, реализующий интерфейс AttributeConverter из состава JPA. При сохранении вызывается метод convertToDatabaseColumn, записывающий порядковый номер статуса в БД. При чтении вызывается convertToEntityAttribute, восстанавливающий элемент перечисления по номеру. Такой подход позволяет хранить статус компактно - одним числом - и при этом работать с ним в коде как с типобезопасным перечислением.

  \par \redline Управление переходами между статусами реализовано в виде конечного автомата непосредственно в классе RideStatus. Каждый элемент перечисления хранит список разрешённых целевых состояний, в которые из него можно перейти. Метод isTransitionAllowed принимает целевой статус и проверяет его наличие в этом списке. Например, из CREATED разрешены переходы только в WAITING\_FOR\_DRIVER и CANCELLED; из COMPLETED - только в PAYMENT\_PENDING. Метод isUpdatable определяет, можно ли изменять данные поездки в текущем статусе. Обновление адресов и цены допускается только на ранних этапах - CREATED, WAITING\_FOR\_DRIVER и DRIVER\_ASSIGNED, - поскольку после того как водитель принял заказ, менять условия поездки уже нельзя. Диаграмма состояний с указанием всех переходов была приведена в главе 3.

  \par \redline API сервиса поездок предоставляет шесть основных методов, перечисленных в таблице \thechaptercntr .\thetablecntr .

  \topTablespace
  {\begin{Center}
      \par Таблица \thechaptercntr .\thetablecntr \spc {--} API сервиса поездок

      \begin{tabular}{|c|c|c|}
        \hline
        Метод & Путь & Назначение \\ \hline
        GET & /api/v1/rides/\{id\} & Получение поездки по идентификатору \\ \hline
        GET & /api/v1/rides & Список всех поездок с пагинацией \\ \hline
        GET & /api/v1/rides/driver/\{driverId\} & Поездки конкретного водителя \\ \hline
        GET & /api/v1/rides/passenger/\{passengerId\} & Поездки конкретного пассажира \\ \hline
        POST & /api/v1/rides & Создание новой поездки \\ \hline
        PATCH & /api/v1/rides/\{id\}/assign-driver & Назначение водителя \\ \hline
        PUT & /api/v1/rides/\{id\} & Изменение данных поездки \\ \hline
        PATCH & /api/v1/rides/\{id\} & Изменение статуса поездки \\ \hline
    \end{tabular} \end{Center}}
  \botTablespace \addtocounter{tablecntr}{1}

  \par \redline Создание поездки - наиболее сложная операция сервиса. Она реализована в методе createRide сервисного слоя и выполняется в несколько строго определённых шагов. Первым делом через Feign-клиент PassengerFeignClient проверяется, что пассажир с заданным идентификатором существует и не занят в другой активной поездке. Если пассажир занят, выбрасывается исключение, и заказ не создаётся - это предотвращает ситуацию, когда один человек создаёт несколько параллельных заказов. Затем проверяется, что адреса отправления и назначения не совпадают, иначе поездка теряет смысл.

  \par \redline После успешной валидации начинается работа с геоданными. Каждый адрес, переданный клиентом в текстовом виде, направляется в картографический сервис. Сервис возвращает координаты, которые сохраняются в таблицу address. Если такой адрес уже существует в БД, новая запись не создаётся - используется существующая. По полученным координатам отправления и назначения через тот же картографический сервис вычисляется расстояние в метрах. Итоговая стоимость рассчитывается по формуле (\thechaptercntr .\theformulacntr):

  \formulaspace \par \redline
    $P = D \cdot T / 1000$,
    \hfill (\thechaptercntr .\theformulacntr) \redline
  \formulaspace \addtocounter{formulacntr}{1}

  \begin{tabular}{p{0,875cm}p{0,3cm}p{15,175cm}}
    & где  & $P$ {--} стоимость поездки в условных единицах; \\
    &      & $D$ {--} дистанция поездки в метрах; \\
    &      & $T$ {--} тариф, равный 10 условных единиц за километр. \\
  \end{tabular}

  \par \redline После расчёта цены создаётся объект Ride. В него записываются ссылки на адреса, дистанция, идентификатор пассажира, рассчитанная цена и текущее время в поле order\_create\_time. Поле ride\_status устанавливается в CREATED. Объект сохраняется в БД через репозиторий. Сразу после сохранения статус переводится в WAITING\_FOR\_DRIVER - заказ готов к тому, чтобы на него назначили водителя. В рамках той же транзакции в таблицу kafka\_messages добавляется запись с топиком passenger-busy-topic и значением true. Это событие позже будет доставлено в сервис пассажиров, и пассажир будет отмечен как занятый.

  \par \redline Назначение водителя выполняется отдельным методом assignDriver. На вход подаются идентификатор поездки и идентификатор водителя. Метод извлекает поездку из БД и проверяет, что она находится в статусе WAITING\_FOR\_DRIVER - назначать водителя на поездку, которая уже выполняется или завершена, бессмысленно и запрещено. Затем через DriverFeignClient запрашивается информация о водителе. Проверяются три условия: водитель должен быть активен (поле is\_enabled = true), не занят в другой поездке (is\_busy = false) и иметь хотя бы один зарегистрированный автомобиль. Если хотя бы одно условие не выполнено, выбрасывается исключение. При успешной проверке в запись Ride записывается driver\_id, статус меняется на DRIVER\_ASSIGNED, а в kafka\_messages добавляется событие driver-busy-topic со значением true.

  \par \redline Изменение статуса поездки выполняется методом updateRideStatus. Он принимает идентификатор поездки и целевой статус. Поездка извлекается из БД, после чего вызывается isTransitionAllowed для проверки допустимости перехода. Если переход разрешён, статус обновляется. При определённых переходах возникают побочные эффекты, реализованные в виде отдельных действий в том же методе.

  \par \redline При переходе в PASSENGER\_PICKED\_UP в поле pickup\_time записывается текущее время - это отметка о моменте, когда водитель забрал пассажира. При переходе в COMPLETED заполняется complete\_time, а в kafka\_messages добавляются два события с топиками driver-busy-topic и passenger-busy-topic со значением false: и водитель, и пассажир освобождаются и могут участвовать в новых поездках. При переходе в RATE в kafka\_messages добавляется событие ride-completed-topic, содержащее идентификаторы поездки, водителя и пассажира, - оно будет доставлено в сервис рейтинга для создания записи ride\_info. При переходе в CANCELLED также публикуются события освобождения участников, если они были назначены к моменту отмены. Без этого шага отменённая поездка заблокировала бы водителя и пассажира навсегда.

  \par \redline Для взаимодействия с другими сервисами через HTTP используются Feign-клиенты из состава Spring Cloud OpenFeign. Сервис поездок содержит два таких клиента: DriverFe-\\ignClient и PassengerFeignClient. Каждый из них описан как интерфейс с аннотацией \\@FeignClient, в которой указано имя сервиса, зарегистрированное в Eureka, - driver-service и passenger-service соответственно. Методы интерфейса аннотированы стандартными аннотациями Spring MVC: @GetMapping, @PathVariable и другими. При вызове метода Feign автоматически строит HTTP-запрос к нужному сервису, сериализует параметры и десериализует ответ. Благодаря интеграции с Eureka не требуется указывать конкретный адрес - Feign получает список доступных экземпляров сервиса от Eureka и распределяет запросы между ними.

  \par \redline Все Feign-клиенты сконфигурированы с использованием библиотеки Resilience4J. Для каждого клиента заданы две политики. Первая - повторные попытки: если вызов завершился ошибкой, предпринимается ещё несколько попыток с возрастающей задержкой между ними. Вторая - circuit breaker: если доля ошибочных вызовов превышает заданный порог, цепь размыкается, и последующие вызовы немедленно возвращают ошибку без фактического обращения к целевому сервису. Через некоторое время цепь переходит в полуоткрытое состояние, и несколько пробных запросов проверяют, восстановилась ли работоспособность. Такая схема предотвращает каскадные отказы: временная недоступность одного сервиса не приводит к исчерпанию ресурсов у вызывающего [1,3].

  \par \redline На уровне контроллера все входящие данные проходят валидацию. Для этого используются аннотации из пакета jakarta.validation: @NotNull, @NotBlank, @Positive и другие. Если данные не проходят проверку, Spring автоматически возвращает клиенту ответ с кодом 400 и описанием ошибок. Для обработки исключений бизнес-логики в сервисе реализован глобальный обработчик с аннотацией @ControllerAdvice. Каждое исключение преобразуется в стандартизованный JSON-ответ с кодом ошибки, текстовым описанием и временной меткой. Такой подход делает API предсказуемым для клиентов и упрощает отладку.

  \par \redline Помимо сервиса поездок, в системе реализованы сервисы водителей и пассажиров. Они устроены значительно проще, поскольку выполняют в основном CRUD-операции над своими сущностями. Сервис водителей работает с двумя сущностями: Driver и Car, связанными отношением "один ко многим". Сервис пассажиров оперирует единственной сущностью Passenger. Оба сервиса предоставляют стандартные методы создания, чтения, обновления и удаления записей, а также метод /me, возвращающий профиль текущего пользователя по идентификатору из заголовка X-User-Id. Удаление реализовано как мягкое: запись не удаляется физически, а помечается как неактивная установкой флага is\_enabled в false. Это позволяет сохранить историю поездок и избежать нарушения ссылочной целостности между сервисами [4].

  \par \redline Сервис рейтинга отвечает за хранение и обработку оценок. Он содержит сущности RideInfo и Rating, связанные отношением "один ко многим". RideInfo создаётся при получении события о завершении поездки из топика ride-completed-topic. После этого водитель и пассажир могут через POST-запрос добавить оценку - числовое значение от 0 до 5 и текстовый комментарий. API сервиса также предоставляет методы для получения среднего рейтинга водителя или пассажира. Средний рейтинг вычисляется динамически: из БД выбираются последние 20 оценок для заданного участника и рассчитывается среднее арифметическое.

  \par
}

\subtitlespace

\subsection*{
  \gostTitleFont
  \redline
  \thechaptercntr .\thesubchaptercntr \spc
  Разработка модуля обработки геолокации и поиска ближайших водителей
} \addtocounter{subchaptercntr}{1}

\subtitlespace

{\gostFont

  \par \redline Работа с геоданными - одна из ключевых особенностей сервиса такси. Пользователь указывает адреса в привычном текстовом виде, однако системе для расчёта стоимости, построения маршрута и передачи информации водителю необходимы географические координаты. Эту задачу решает модуль геолокации, взаимодействующий с картографическим сервисом и обеспечивающий хранение и кэширование полученных данных [6].

  \par \redline Взаимодействие с сервисом Mapbox инкапсулировано в отдельном классе \\MapboxClient. Он содержит два публичных метода. Первый метод - geocode - принимает название адреса строкой и возвращает координаты: широту (latitude) и долготу (longitude) в виде чисел типа Double. Второй метод - getDistance - принимает две пары координат и возвращает расстояние между ними в метрах, также в виде Double. HTTP-запросы к Mapbox выполняются с помощью RestClient из состава Spring Boot [5].

  \par \redline Для работы с Mapbox требуется API-ключ, который передаётся в заголовке каждого запроса. Ключ хранится в конфигурации сервиса и не попадает в исходный код. В конфигурационном файле заданы базовые URL API геокодирования и API маршрутизации, а также параметры по умолчанию, такие как ограничение области поиска [6].

  \par \redline При вызове geocode формируется GET-запрос к API геокодирования Mapbox. В URL передаётся параметр поиска, содержащий название адреса. Сервис возвращает массив совпадений в формате JSON, из которого клиентский код извлекает координаты первого - наиболее релевантного - результата. Если сервис не находит ни одного совпадения, генерируется исключение, и создание заказа прерывается.

  \par \redline Метод getDistance формирует запрос к API маршрутизации Mapbox. В запросе передаются две точки с координатами, а также параметр, указывающий, что требуется автомобильный маршрут. Сервис возвращает JSON-ответ, содержащий секцию с рассчитанной дистанцией в метрах - именно это значение извлекается и возвращается клиентским кодом.

  \par \redline Следует пояснить, почему для расчёта расстояния используется внешний сервис, а не простая математическая формула. Расстояние между двумя точками на поверхности Земли можно вычислить по формуле гаверсинусов, которая определяет кратчайшее расстояние по дуге большого круга:

  \formulaspace \par \redline
    $d = 2 \cdot R \cdot \arcsin\left(\sqrt{\sin^2\left(\frac{\Delta\phi}{2}\right) +
    \cos(\phi_1) \cdot \cos(\phi_2) \cdot \sin^2\left(\frac{\Delta\lambda}{2}\right)}\right)$,
    \hfill (\thechaptercntr .\theformulacntr) \redline
  \formulaspace \addtocounter{formulacntr}{1}

  \begin{tabular}{p{0,875cm}p{0,3cm}p{15,175cm}}
    & где  & $R$ {--} средний радиус Земли, равный 6371000 м; \\
    &      & $\phi_1, \phi_2$ {--} широты первой и второй точек в радианах; \\
    &      & $\lambda_1, \lambda_2$ {--} долготы первой и второй точек в радианах; \\
    &      & $\Delta\phi = \phi_2 - \phi_1$; \\
    &      & $\Delta\lambda = \lambda_2 - \lambda_1$. \\
  \end{tabular}

  \par \redline Формула (4.2) даёт расстояние по прямой. Однако в приложении такси пассажира интересует не это расстояние, а фактический путь по дорогам. Именно дорожная дистанция определяет длительность поездки и её стоимость. Mapbox Directions API учитывает дорожную сеть, разрешённые направления движения, скоростные ограничения и даже рельеф местности, возвращая реальную дистанцию по автомобильным дорогам. Разница между гаверсинусами и дорожной дистанцией может составлять десятки процентов, особенно в городах с плотной застройкой и сложной дорожной сетью. Поэтому в проекте используется именно Mapbox API [6].

  \par \redline Запросы к картографическому сервису выполняются по сети и имеют заметную задержку - от десятков до сотен миллисекунд. При этом результаты для одних и тех же адресов не меняются между вызовами: координаты здания на карте постоянны, а расстояние по дорогам между двумя точками меняется редко. Поэтому в системе реализовано трёхуровневое кэширование с использованием Redis [7].

  \par \redline Первый уровень - кэширование самого адреса как сущности. Когда сервис получает название адреса, он сначала ищет его в таблице address по полю address\_name. Если запись найдена, она возвращается немедленно, без запроса к картографическому сервису. Для быстрого поиска по названию на поле address\_name создан индекс. Второй уровень - кэширование результата геокодирования. Даже если адрес ещё не сохранён в БД, координаты, полученные для конкретного названия, помещаются в Redis с ограниченным временем жизни. При повторном запросе того же названия координаты достаются из кэша. Третий уровень - кэширование дистанции. Расстояние между конкретными координатами сохраняется в Redis, и при повторном расчёте для той же пары точек сетевой запрос не выполняется.

  \par \redline Технически кэширование реализовано с помощью механизма аннотаций Spring Cache. В конфигурационном классе сервиса поездок определён менеджер кэша, использующий Redis в качестве хранилища. Для каждого уровня кэширования заведён свой кэш с уникальным именем: addresses, geocode и distance. Методы сервисного слоя, отвечающие за получение адресов, координат и дистанций, помечены аннотацией @Cacheable с указанием имени соответствующего кэша. Ключ кэша формируется автоматически на основе значений параметров метода. При вызове такого метода Spring сначала проверяет наличие значения в Redis. Если значение найдено - оно возвращается, минуя тело метода. Если нет - метод выполняется, а его результат помещается в Redis. Время жизни записей задано разным для разных кэшей: адреса хранятся дольше, результаты геокодирования - меньше, а дистанции имеют самое короткое время жизни, поскольку дорожная ситуация может измениться [7].

  \par \redline Координаты, полученные от картографического сервиса и сохранённые в таблице address, имеют тип double precision. Широта принимает значения в диапазоне от -90 до 90 градусов, долгота - от -180 до 180. Перед сохранением оба значения проверяются на принадлежность этим диапазонам. Проверка выполняется на уровне приложения, в сервисном слое, до передачи данных в репозиторий. Если координаты выходят за допустимые границы, операция прерывается с исключением. Это исключает сохранение заведомо некорректных данных, которые могли бы возникнуть, например, при ошибке картографического сервиса.

  \par \redline В текущей версии системы поиск ближайших водителей также реализован через картографический сервис. При запросе на назначение водителя сервис поездок получает список свободных водителей и передаёт их координаты вместе с координатами пассажира в картографический сервис, который упорядочивает водителей по расстоянию. Такой подход не требует пространственных запросов к БД и использования расширений вроде PostGIS. Однако модель данных спроектирована так, что в будущем можно добавить к таблице address столбец геометрического типа и заполнить его значениями из уже имеющихся полей широты и долготы. Это потребует минимальных изменений в коде и схеме БД [6].

  \par
}

\subtitlespace

\subsection*{
  \gostTitleFont
  \redline
  \thechaptercntr .\thesubchaptercntr \spc
  Система авторизации и аутентификации на основе JWT
} \addtocounter{subchaptercntr}{1}

\subtitlespace

{\gostFont

  \par \redline Вопросы безопасности в сервисе такси стоят особенно остро, поскольку система работает с персональными данными пассажиров и водителей, историей поездок и платёжной информацией. Для обеспечения защиты в проекте используется Keycloak - сервер авторизации с открытым исходным кодом, работающий по протоколу OAuth2 и использующий JWT в качестве формата токенов доступа. Keycloak выполняет роль единого Identity Provider: только он хранит учётные данные пользователей и управляет ролями, а все остальные сервисы лишь проверяют подлинность предъявленных токенов [5].

  \par \redline В Keycloak для приложения создан отдельный realm, изолирующий его настройки от других систем, которые могут работать на том же сервере. В рамках realm определены клиенты для каждого микросервиса, включая API Gateway. Также определены роли, отражающие уровни доступа в системе. Роль USER назначается каждому зарегистрированному пользователю автоматически и предоставляет базовые права на использование системы. Роли DRIVER и PASSENGER указывают на бизнес-роль пользователя и определяют, с каким сервисом - водителей или пассажиров - он работает. Роль ADMIN предназначена для административных операций. Отдельная роль SERVICE используется для межсервисной авторизации: когда один микросервис обращается к другому, он действует от имени этой роли, а не от имени конечного пользователя.

  \par \redline Сервис аутентификации предоставляет клиентам четыре основных метода, перечисленных в таблице \thechaptercntr .\thetablecntr .

  \topTablespace
  {\begin{Center}
      \par Таблица \thechaptercntr .\thetablecntr \spc {--} API сервиса аутентификации

      \begin{tabular}{|c|c|c|}
        \hline
        Метод & Путь & Назначение \\ \hline
        POST & /api/v1/auth/signup & Регистрация пользователя \\ \hline
        POST & /api/v1/auth/signin & Вход и получение токенов \\ \hline
        POST & /api/v1/auth/refresh & Обновление access-токена \\ \hline
        POST & /api/v1/auth/logout & Отзыв refresh-токена \\ \hline
    \end{tabular} \end{Center}}
  \botTablespace \addtocounter{tablecntr}{1}

  \par \redline Процесс регистрации нового пользователя включает взаимодействие трёх компонентов: сервиса аутентификации, Keycloak и одного из бизнес-сервисов. Клиент отправляет POST-запрос на /api/v1/auth/signup, передавая в теле имя, фамилию, адрес электронной почты, пароль, номер телефона и желаемую бизнес-роль - DRIVER или PASSENGER. Сервис аутентификации принимает запрос, валидирует входные данные и приступает к созданию учётной записи.

  \par \redline На первом шаге через Keycloak Admin Client создаётся пользователь в Keycloak. Для этого используется специальная учётная запись администратора realm, настроенная в конфигурации сервиса. Keycloak Admin Client принимает данные пользователя и создаёт запись во внутреннем хранилище Keycloak, возвращая уникальный идентификатор нового пользователя. На втором шаге пользователю назначаются роли: обязательно роль USER, а также выбранная бизнес-роль - DRIVER или PASSENGER. Назначение ролей также выполняется через Keycloak Admin Client. На третьем шаге сервис аутентификации через HTTP-запрос обращается к сервису водителей или пассажиров - в зависимости от выбранной роли - и создаёт там профиль пользователя, передавая его идентификатор из Keycloak, имя, фамилию, email и телефон. Все три шага образуют единый бизнес-процесс регистрации, и если на любом из них происходит ошибка, процесс прерывается.

  \par \redline Для входа в систему клиент отправляет POST /signin с email и паролем. Сервис аутентификации направляет эти учётные данные в Keycloak, используя стандартный механизм OAuth2 - Password Grant Flow. В запросе к Keycloak передаются client\_id, client\_secret, имя пользователя, пароль и scope. Keycloak проверяет учётные данные и, если они корректны, формирует и возвращает два JWT-токена: access-токен и refresh-токен. Access-токен имеет короткое время жизни и содержит в себе идентификатор пользователя в поле sub, его роли в поле realm\_access, а также имя клиента, для которого выдан токен. Refresh-токен имеет более длительное время жизни и используется только для получения нового access-токена, когда срок действия текущего истёк. Оба токена возвращаются клиенту в JSON-ответе.

  \par \redline Для обновления access-токена клиент отправляет POST /refresh, передавая refresh-токен. Сервис аутентификации перенаправляет запрос в Keycloak, который проверяет \\refresh-токен и, если он действителен, выпускает новую пару токенов. Для завершения сессии клиент отправляет POST /logout с refresh-токеном, и Keycloak отзывает его, делая невозможным дальнейшее обновление access-токена.

  \par \redline Проверка токенов при каждом запросе выполняется не сервисом аутентификации, а непосредственно на уровне API Gateway. Шлюз сконфигурирован как OAuth2 Resource Server. В его настройках указан адрес Keycloak и публичный ключ для проверки подписи JWT. При получении запроса шлюз извлекает токен из заголовка Authorization, проверяет его подпись, срок действия и соответствие эмитента. Если проверка не пройдена, шлюз возвращает ответ с кодом 401, и запрос не доходит до бизнес-сервисов. Все пути, начинающиеся с /api/v1/auth и /actuator, исключены из проверки - они должны быть доступны без авторизации [5].

  \par \redline После успешной проверки шлюз выполняет дополнительный шаг: он извлекает из токена идентификатор пользователя из поля sub и помещает его в HTTP-заголовок X-User-Id. Этот заголовок добавляется к запросу при его перенаправлении к целевому сервису. Реализовано это в фильтре AddUserIdHeaderFilter, который является компонентом Spring Cloud Gateway. Фильтр выполняется после проверки токена и перед отправкой запроса нижестоящему сервису. Благодаря этому заголовку любой бизнес-сервис может определить, от имени какого пользователя пришёл запрос.

  \par \redline Наличие заголовка X-User-Id в каждом запросе позволило реализовать в сервисах водителей и пассажиров удобный адрес /me. Клиенту не нужно знать внутренний идентификатор пользователя в системе - достаточно обратиться к /me, и сервис, прочитав X-User-Id из заголовка, найдёт соответствующий профиль и вернёт его. Это упрощает клиентский код и повышает безопасность, поскольку идентификаторы не передаются в URL [4].

  \par \redline Межсервисное взаимодействие защищено отдельным механизмом. Когда один сервис обращается к другому через Feign-клиент, он не использует токен конечного пользователя. Вместо этого применяется Client Credentials Flow - ещё один стандартный механизм OAuth2, предназначенный для взаимодействия серверов без участия пользователя. Сервис отправляет в Keycloak свои client\_id и client\_secret и получает служебный access-токен, ассоциированный с ролью SERVICE. Этот токен передаётся в заголовке Authorization Feign-запроса. Целевой сервис проверяет токен так же, как и пользовательский, но видит в нём роль SERVICE и предоставляет доступ в соответствии с этой ролью. Таким образом достигается баланс между удобством межсервисного взаимодействия и контролем доступа.

  \par
}

\subtitlespace

\subsection*{
  \gostTitleFont
  \redline
  \thechaptercntr .\thesubchaptercntr \spc
  Интеграция брокера сообщений для обработки событий в реальном времени
} \addtocounter{subchaptercntr}{1}

\subtitlespace

{\gostFont

  \par \redline В системе заказа такси многие действия порождают события, о которых должны быть уведомлены другие сервисы. При создании поездки пассажир переходит в состояние занятости. При назначении водителя он также помечается занятым. При завершении или отмене поездки оба участника освобождаются, а сервис рейтинга получает информацию для создания записи о выполненном заказе. Для передачи таких событий в проекте используется брокер сообщений Apache Kafka - распределённая платформа потоковой обработки, обеспечивающая надёжную асинхронную доставку сообщений между компонентами системы [8].

  \par \redline Kafka была выбрана по нескольким причинам. Во-первых, она обеспечивает сохранение сообщений на диске и их репликацию между брокерами, что гарантирует доставку даже при временных сбоях. Во-вторых, Kafka поддерживает модель издатель-подписчик, при которой одно сообщение может быть доставлено нескольким независимым получателям. В-третьих, Kafka хранит сообщения упорядоченно в рамках одного раздела топика, что позволяет восстанавливать последовательность событий при необходимости.

  \par \redline В системе определены три топика, сведённые в таблицу \thechaptercntr .\thetablecntr .
  \topTablespace
  {\begin{Center}
      \par Таблица \thechaptercntr .\thetablecntr \spc {--} Топики Kafka

      \begin{tabular}{|c|c|c|p{4cm}|}
        \hline
        Топик & Отправитель & Получатель & Назначение \\ \hline
        passenger-busy-topic & сервис поездок & сервис пассажиров & Изменение статуса занятости \\ \hline
        driver-busy-topic & сервис поездок & сервис водителей & Изменение статуса занятости \\ \hline
        ride-completed-topic & сервис поездок & сервис рейтинга & Уведомление о завершении поездки \\ \hline
    \end{tabular} \end{Center}}
  \botTablespace \addtocounter{tablecntr}{1}

  \par \redline Топики passenger-busy-topic и driver-busy-topic устроены идентично. Ключом сообщения выступает идентификатор пассажира или водителя (Long), а телом - булево значение, инкапсулированное в JSON. Если значение равно true, участник помечается как занятый; если false - как свободный. Использование идентификатора в качестве ключа гарантирует, что все сообщения, относящиеся к одному участнику, попадают в один и тот же раздел топика и обрабатываются в порядке поступления - это важно, чтобы статусы занятости не перепутались.

  \par \redline Топик ride-completed-topic использует идентификатор поездки в качестве ключа. Телом сообщения служит JSON-объект с тремя полями: rideId, driverId и passengerId. Этих данных достаточно сервису рейтинга, чтобы создать запись ride\_info, к которой впоследствии будут привязаны оценки.

  \par \redline Отправка сообщений реализована с применением паттерна Outbox. Его идея состоит в том, что сообщение не отправляется в Kafka непосредственно из сервисного слоя после изменения данных. Вместо этого оно сначала записывается в таблицу kafka\_messages - в той же транзакции БД, что и изменение данных поездки. Транзакция либо фиксируется целиком - и тогда и новые данные, и сообщение гарантированно сохранены, - либо откатывается, и тогда не сохраняется ни то, ни другое. Промежуточных состояний не возникает. Такой подход решает классическую проблему распределённых систем: когда данные уже изменены, а событие ещё не отправлено (или наоборот), что приводит к несогласованности между сервисами. С Outbox такого не происходит [8].

  \par \redline Таблица kafka\_messages спроектирована для хранения всей необходимой информации об исходящем сообщении. Поле message\_id - автоинкрементный первичный ключ. Поле topic содержит имя топика Kafka. Поле key - ключ сообщения типа Long. Поле message хранит тело сообщения в виде строки в формате JSON, для этого использован тип TEXT. Поля created\_at и sent\_at - временные метки, заполняемые при создании и отправке соответственно. Поле is\_sent - булев флаг, по умолчанию равный false и меняющийся на true после успешной отправки. Также в таблице присутствуют поля trace\_id и span\_id, используемые для сквозной распределённой трассировки запросов. На полях is\_sent и created\_at создан составной индекс: он делает эффективным основной запрос планировщика, выбирающий неотправленные сообщения в порядке их создания.

  \par \redline Фактическая отправка сообщений из таблицы в Kafka выполняется фоновым планировщиком KafkaMessageScheduler. Он реализован как метод с аннотацией @Scheduled, настроенной на выполнение с фиксированной задержкой. В каждом цикле планировщик выполняет запрос к репозиторию KafkaMessage, выбирая все записи, у которых is\_sent = false, в порядке возрастания created\_at. Для каждой записи вызывается метод \\KafkaTemplate.send, передающий топик, ключ и тело сообщения. После успешного вызова поле is\_sent устанавливается в true, а в sent\_at записывается текущее время. Если отправка конкретного сообщения завершается ошибкой, оно остаётся в таблице с is\_sent = false и будет повторно обработано в следующем цикле.

  \par \redline Так как сервис поездок может быть запущен в нескольких экземплярах - например, для повышения отказоустойчивости, - планировщик должен работать только на одном из них. Иначе два экземпляра одновременно прочитают одни и те же неотправленные сообщения и опубликуют их дважды. Проблема решается с помощью библиотеки ShedLock, реализующей механизм распределённых блокировок на основе БД. В конфигурации ShedLock задаётся поставщик блокировок, указывающий на основную БД сервиса, и создаётся служебная таблица shedlock. Перед выполнением метода планировщика ShedLock пытается вставить в эту таблицу запись с уникальным именем задачи. Если запись успешно создана - блокировка получена, и планировщик выполняет свою работу. Если запись уже существует - значит, блокировку держит другой экземпляр, и текущий цикл пропускается. По завершении метода блокировка освобождается, либо истекает по таймауту, если экземпляр аварийно завершился [8].

  \par \redline На стороне получателей сообщения обрабатываются с помощью классов, аннотированных @KafkaListener из библиотеки Spring for Apache Kafka. В сервисе пассажиров слушатель подписан на топик passenger-busy-topic. При получении сообщения он извлекает ключ - идентификатор пассажира - и значение занятости из тела, десериализуя JSON в Boolean. Затем через репозиторий находится запись пассажира, и поле is\_busy обновляется в соответствии с полученным значением. Аналогично работает слушатель в сервисе водителей для топика driver-busy-topic. Оба слушателя спроектированы как идемпотентные: повторная обработка одного и того же сообщения не приводит к некорректному состоянию, поскольку каждая обработка просто устанавливает поле is\_busy в переданное значение.

  \par \redline Слушатель в сервисе рейтинга подписан на топик ride-completed-topic. При получении сообщения он десериализует JSON в объект класса RideInfoPayload, содержащий rideId, driverId и passengerId. Затем через репозиторий создаётся новая запись ride\_info с этими идентификаторами. После создания записи и водитель, и пассажир могут выставить оценки через POST /api/v1/ratings, указав идентификатор поездки, числовую оценку, текстовый комментарий и признак того, кто оценивает - пассажир или водитель. Слушатель также является идемпотентным благодаря уникальному ограничению на ride\_id в таблице ride\_info, предотвращающему создание дубликатов [8].

  \par \redline Для предотвращения бесконтрольного роста таблицы kafka\_messages реализован отдельный планировщик очистки. Раз в сутки - в ночное время, когда нагрузка на систему минимальна - он выполняет запрос на удаление всех записей, у которых is\_sent = true и sent\_at старше одного дня. Хранение отправленных сообщений дольше этого срока не требуется, поскольку все заинтересованные сервисы уже получили и обработали их.

  \par \redline Конфигурация Kafka задаётся в файлах настроек Spring Boot каждого сервиса, использующего брокер. В настройках указываются адреса брокеров, идентификатор группы потребителей, стратегия подтверждения доставки и параметры сериализации. Ключи сообщений сериализуются стандартным LongSerializer, значения - StringSerializer. На стороне потребителей применяются соответствующие десериализаторы. Для повышения надёжности все топики сконфигурированы с фактором репликации, обеспечивающим сохранность данных при выходе из строя одного из узлов кластера. Также настроено автоматическое создание топиков при запуске приложения, что упрощает развёртывание системы в новой среде [8].

  \par \redline Особенностью реализации является интеграция с системой распределённой трассировки. При сохранении сообщения в таблицу kafka\_messages в поля trace\_id и span\_id записываются соответствующие идентификаторы из текущего контекста трассировки. При отправке сообщения в Kafka эти идентификаторы восстанавливаются и передаются в заголовках сообщения. На стороне потребителя они снова извлекаются и помещаются в контекст трассировки. Благодаря этому в системах мониторинга можно проследить полный путь запроса - от HTTP-вызова до обработки Kafka-сообщения и обратно.

  \par \redline Таким образом, выбранный подход к организации асинхронного взаимодействия обеспечивает надёжную доставку событий между сервисами без жёстких синхронных зависимостей. Паттерн Outbox гарантирует согласованность данных при любом исходе транзакции, ShedLock предотвращает дублирование сообщений при многократном запуске отправителя, а идемпотентные слушатели делают обработку устойчивой к повторным доставкам.

  \par
}

\setcounter{subchaptercntr}{1}
\setcounter{formulacntr}{1}
\setcounter{imagecntr}{1}
\setcounter{tablecntr}{1}
